\chapter{Unit 1: Introduction, Variables, Expressions and Statements}

\section{Section A: Multiple Choice Practice Questions (60 Curated Problems)}

\noindent This section contains 60 rigorous multiple-choice questions extracted from university course examination assessments. Review each question, test edge-case syntax, and verify code execution mechanics.
\vspace{3mm}

\noindent\textbf{Q1.} \hfill {\small\color{LPUOrange}\textbf{[need for a programming language | Easy]}}\\[2pt]
What is the primary purpose of a programming language?

\begin{enumerate}[label=(\alph*),topsep=2pt,itemsep=1pt,partopsep=0pt,parsep=0pt]
    \item \relax To store data permanently on disk
    \item \relax To increase the speed of the CPU
    \item \relax To design computer hardware circuits
    \item \relax To communicate instructions to a computer
\end{enumerate}
\vspace{2mm}

\noindent\textbf{Q2.} \hfill {\small\color{LPUOrange}\textbf{[introduction to python as a programming language | Easy]}}\\[2pt]
Python is best described as which type of programming language?

\begin{enumerate}[label=(\alph*),topsep=2pt,itemsep=1pt,partopsep=0pt,parsep=0pt]
    \item \relax A low-level assembly language
    \item \relax A markup language for web pages
    \item \relax A high-level, interpreted language
    \item \relax A hardware description language
\end{enumerate}
\vspace{2mm}

\noindent\textbf{Q3.} \hfill {\small\color{LPUOrange}\textbf{[introduction to python as a programming language | Easy]}}\\[2pt]
Who is credited with creating the Python programming language?

\begin{enumerate}[label=(\alph*),topsep=2pt,itemsep=1pt,partopsep=0pt,parsep=0pt]
    \item \relax Dennis Ritchie
    \item \relax James Gosling
    \item \relax Guido van Rossum
    \item \relax Bjarne Stroustrup
\end{enumerate}
\vspace{2mm}

\noindent\textbf{Q4.} \hfill {\small\color{LPUOrange}\textbf{[programming errors and debugging | Easy]}}\\[2pt]
Which type of error occurs when Python code violates the rules of the language, such as a missing colon?

\begin{enumerate}[label=(\alph*),topsep=2pt,itemsep=1pt,partopsep=0pt,parsep=0pt]
    \item \relax Runtime error
    \item \relax Semantic error
    \item \relax Logical error
    \item \relax Syntax error
\end{enumerate}
\vspace{2mm}

\noindent\textbf{Q5.} \hfill {\small\color{LPUOrange}\textbf{[programming errors and debugging | Easy]}}\\[2pt]
What is the process of finding and fixing errors in a program called?

\begin{enumerate}[label=(\alph*),topsep=2pt,itemsep=1pt,partopsep=0pt,parsep=0pt]
    \item \relax Debugging
    \item \relax Compiling
    \item \relax Executing
    \item \relax Formatting
\end{enumerate}
\vspace{2mm}

\noindent\textbf{Q6.} \hfill {\small\color{LPUOrange}\textbf{[identifiers | Easy]}}\\[2pt]
Which of the following is a valid Python identifier?

\begin{enumerate}[label=(\alph*),topsep=2pt,itemsep=1pt,partopsep=0pt,parsep=0pt]
    \item \relax class
    \item \relax 2ndValue
    \item \relax \_myVar
    \item \relax my-var
\end{enumerate}
\vspace{2mm}

\noindent\textbf{Q7.} \hfill {\small\color{LPUOrange}\textbf{[identifiers | Easy]}}\\[2pt]
Which statement about Python identifiers is true?

\begin{enumerate}[label=(\alph*),topsep=2pt,itemsep=1pt,partopsep=0pt,parsep=0pt]
    \item \relax They are case-sensitive
    \item \relax They may contain spaces
    \item \relax They ignore uppercase letters
    \item \relax They can start with a digit
\end{enumerate}
\vspace{2mm}

\noindent\textbf{Q8.} \hfill {\small\color{LPUOrange}\textbf{[variables | Easy]}}\\[2pt]
What is a variable in Python?

\begin{enumerate}[label=(\alph*),topsep=2pt,itemsep=1pt,partopsep=0pt,parsep=0pt]
    \item \relax A fixed value that never changes
    \item \relax A built-in function name
    \item \relax A type of loop structure
    \item \relax A named reference to a stored value
\end{enumerate}
\vspace{2mm}

\noindent\textbf{Q9.} \hfill {\small\color{LPUOrange}\textbf{[assignment statements | Easy]}}\\[2pt]
What does the statement \texttt{x = 5} do in Python?

\begin{enumerate}[label=(\alph*),topsep=2pt,itemsep=1pt,partopsep=0pt,parsep=0pt]
    \item \relax Prints the value 5
    \item \relax Checks if x equals 5
    \item \relax Assigns the value 5 to the variable x
    \item \relax Compares x with 5
\end{enumerate}
\vspace{2mm}

\noindent\textbf{Q10.} \hfill {\small\color{LPUOrange}\textbf{[expressions | Easy]}}\\[2pt]
Which of the following is an example of an expression in Python?

\begin{enumerate}[label=(\alph*),topsep=2pt,itemsep=1pt,partopsep=0pt,parsep=0pt]
    \item \relax 3 +
    \item \relax while True:
    \item \relax if x:
    \item \relax def foo():
\end{enumerate}
\vspace{2mm}

\noindent\textbf{Q11.} \hfill {\small\color{LPUOrange}\textbf{[named constant | Easy]}}\\[2pt]
By convention, how are named constants usually written in Python?

\begin{enumerate}[label=(\alph*),topsep=2pt,itemsep=1pt,partopsep=0pt,parsep=0pt]
    \item \relax In all lowercase letters
    \item \relax In all uppercase letters
    \item \relax In camelCase
    \item \relax With a leading digit
\end{enumerate}
\vspace{2mm}

\noindent\textbf{Q12.} \hfill {\small\color{LPUOrange}\textbf{[simultaneous assignment | Easy]}}\\[2pt]
What will be the values of \texttt{a} and \texttt{b} after \texttt{a, b = 1, 2}?

\begin{enumerate}[label=(\alph*),topsep=2pt,itemsep=1pt,partopsep=0pt,parsep=0pt]
    \item \relax both a and b are 2
    \item \relax both a and b are 1
    \item \relax a is 2 and b is 1
    \item \relax a is 1 and b is 2
\end{enumerate}
\vspace{2mm}

\noindent\textbf{Q13.} \hfill {\small\color{LPUOrange}\textbf{[simultaneous assignment | Easy]}}\\[2pt]
Which single statement correctly swaps the values of \texttt{x} and \texttt{y} in Python?

\begin{enumerate}[label=(\alph*),topsep=2pt,itemsep=1pt,partopsep=0pt,parsep=0pt]
    \item \relax x == y, y == x
    \item \relax swap(x, y)
    \item \relax x, y = y, x
    \item \relax x = y = x
\end{enumerate}
\vspace{2mm}

\noindent\textbf{Q14.} \hfill {\small\color{LPUOrange}\textbf{[boolean types | Easy]}}\\[2pt]
Which of the following are the two Boolean values in Python?

\begin{enumerate}[label=(\alph*),topsep=2pt,itemsep=1pt,partopsep=0pt,parsep=0pt]
    \item \relax true and false
    \item \relax True and False
    \item \relax 1 and 0 only
    \item \relax Yes and No
\end{enumerate}
\vspace{2mm}

\noindent\textbf{Q15.} \hfill {\small\color{LPUOrange}\textbf{[numeric data types | Easy]}}\\[2pt]
Which numeric data type represents whole numbers in Python?

\begin{enumerate}[label=(\alph*),topsep=2pt,itemsep=1pt,partopsep=0pt,parsep=0pt]
    \item \relax complex
    \item \relax str
    \item \relax float
    \item \relax int
\end{enumerate}
\vspace{2mm}

\noindent\textbf{Q16.} \hfill {\small\color{LPUOrange}\textbf{[numeric data types | Easy]}}\\[2pt]
What is the data type of the value \texttt{3.14} in Python?

\begin{enumerate}[label=(\alph*),topsep=2pt,itemsep=1pt,partopsep=0pt,parsep=0pt]
    \item \relax complex
    \item \relax float
    \item \relax bool
    \item \relax int
\end{enumerate}
\vspace{2mm}

\noindent\textbf{Q17.} \hfill {\small\color{LPUOrange}\textbf{[operators | Easy]}}\\[2pt]
Which operator is used for exponentiation in Python?

\begin{enumerate}[label=(\alph*),topsep=2pt,itemsep=1pt,partopsep=0pt,parsep=0pt]
    \item \relax //
    \item \relax \textasciicircum{}
    \item \relax \%
    \item \relax **
\end{enumerate}
\vspace{2mm}

\noindent\textbf{Q18.} \hfill {\small\color{LPUOrange}\textbf{[operator precedence and associativity | Easy]}}\\[2pt]
What is the result of the expression \texttt{2 + 3 * 4} in Python?

\begin{enumerate}[label=(\alph*),topsep=2pt,itemsep=1pt,partopsep=0pt,parsep=0pt]
    \item \relax 14
    \item \relax 20
    \item \relax 24
    \item \relax 11
\end{enumerate}
\vspace{2mm}

\noindent\textbf{Q19.} \hfill {\small\color{LPUOrange}\textbf{[augmented assignment operators | Easy]}}\\[2pt]
What does the statement \texttt{x += 3} mean in Python?

\begin{enumerate}[label=(\alph*),topsep=2pt,itemsep=1pt,partopsep=0pt,parsep=0pt]
    \item \relax x = 3
    \item \relax x = x + 3
    \item \relax x = x * 3
    \item \relax x = 3 - x
\end{enumerate}
\vspace{2mm}

\noindent\textbf{Q20.} \hfill {\small\color{LPUOrange}\textbf{[type conversion and rounding | Easy]}}\\[2pt]
What does the function \texttt{int(4.9)} return in Python?

\begin{enumerate}[label=(\alph*),topsep=2pt,itemsep=1pt,partopsep=0pt,parsep=0pt]
    \item \relax 4
    \item \relax 5
    \item \relax 4.9
    \item \relax 0
\end{enumerate}
\vspace{2mm}

\noindent\textbf{Q21.} \hfill {\small\color{LPUOrange}\textbf{[numeric data types | Medium]}}\\[2pt]
What is the data type of the result of the expression \texttt{7 / 2} in Python 3?

\begin{enumerate}[label=(\alph*),topsep=2pt,itemsep=1pt,partopsep=0pt,parsep=0pt]
    \item \relax float
    \item \relax int
    \item \relax decimal
    \item \relax complex
\end{enumerate}
\vspace{2mm}

\noindent\textbf{Q22.} \hfill {\small\color{LPUOrange}\textbf{[operators | Medium]}}\\[2pt]
What is the output of \texttt{17 \% 5} in Python?

\begin{enumerate}[label=(\alph*),topsep=2pt,itemsep=1pt,partopsep=0pt,parsep=0pt]
    \item \relax 2
    \item \relax 12
    \item \relax 3
    \item \relax 5
\end{enumerate}
\vspace{2mm}

\noindent\textbf{Q23.} \hfill {\small\color{LPUOrange}\textbf{[operator precedence and associativity | Medium]}}\\[2pt]
What is the value of the expression \texttt{2 + 3 * 4 ** 2}?

\begin{enumerate}[label=(\alph*),topsep=2pt,itemsep=1pt,partopsep=0pt,parsep=0pt]
    \item \relax 400
    \item \relax 50
    \item \relax 80
    \item \relax 196
\end{enumerate}
\vspace{2mm}

\noindent\textbf{Q24.} \hfill {\small\color{LPUOrange}\textbf{[operator precedence and associativity | Medium]}}\\[2pt]
What does the expression \texttt{2 ** 3 ** 2} evaluate to in Python?

\begin{enumerate}[label=(\alph*),topsep=2pt,itemsep=1pt,partopsep=0pt,parsep=0pt]
    \item \relax 512
    \item \relax 256
    \item \relax 64
    \item \relax 128
\end{enumerate}
\vspace{2mm}

\noindent\textbf{Q25.} \hfill {\small\color{LPUOrange}\textbf{[augmented assignment operators | Medium]}}\\[2pt]
After the code \texttt{x = 10} followed by \texttt{x //= 3}, what is the value of \texttt{x}?

\begin{enumerate}[label=(\alph*),topsep=2pt,itemsep=1pt,partopsep=0pt,parsep=0pt]
    \item \relax 3.33
    \item \relax 3
    \item \relax 4
    \item \relax 1
\end{enumerate}
\vspace{2mm}

\noindent\textbf{Q26.} \hfill {\small\color{LPUOrange}\textbf{[simultaneous assignment | Medium]}}\\[2pt]
After executing \texttt{a, b = 5, 10} and then \texttt{a, b = b, a}, what are the values of \texttt{a} and \texttt{b}?

\begin{enumerate}[label=(\alph*),topsep=2pt,itemsep=1pt,partopsep=0pt,parsep=0pt]
    \item \relax a = 10, b = 5
    \item \relax a = 5, b = 5
    \item \relax a = 10, b = 10
    \item \relax a = 5, b = 10
\end{enumerate}
\vspace{2mm}

\noindent\textbf{Q27.} \hfill {\small\color{LPUOrange}\textbf{[type conversion and rounding | Medium]}}\\[2pt]
What is the output of \texttt{int(9.99)} in Python?

\begin{enumerate}[label=(\alph*),topsep=2pt,itemsep=1pt,partopsep=0pt,parsep=0pt]
    \item \relax 10
    \item \relax 9.99
    \item \relax 9
    \item \relax 9.0
\end{enumerate}
\vspace{2mm}

\noindent\textbf{Q28.} \hfill {\small\color{LPUOrange}\textbf{[type conversion and rounding | Medium]}}\\[2pt]
What does \texttt{round(2.5)} return in Python 3?

\begin{enumerate}[label=(\alph*),topsep=2pt,itemsep=1pt,partopsep=0pt,parsep=0pt]
    \item \relax 2
    \item \relax 2.0
    \item \relax 2.5
    \item \relax 3
\end{enumerate}
\vspace{2mm}

\noindent\textbf{Q29.} \hfill {\small\color{LPUOrange}\textbf{[boolean types | Medium]}}\\[2pt]
What is the result of the expression \texttt{True + True + False} in Python?

\begin{enumerate}[label=(\alph*),topsep=2pt,itemsep=1pt,partopsep=0pt,parsep=0pt]
    \item \relax 2
    \item \relax TypeError
    \item \relax True
    \item \relax 1
\end{enumerate}
\vspace{2mm}

\noindent\textbf{Q30.} \hfill {\small\color{LPUOrange}\textbf{[boolean types | Medium]}}\\[2pt]
What is the output of \texttt{bool(0.0) or bool("")}?

\begin{enumerate}[label=(\alph*),topsep=2pt,itemsep=1pt,partopsep=0pt,parsep=0pt]
    \item \relax True
    \item \relax False
    \item \relax 0.0
    \item \relax None
\end{enumerate}
\vspace{2mm}

\noindent\textbf{Q31.} \hfill {\small\color{LPUOrange}\textbf{[identifiers | Medium]}}\\[2pt]
Which of the following is a valid Python identifier?

\begin{enumerate}[label=(\alph*),topsep=2pt,itemsep=1pt,partopsep=0pt,parsep=0pt]
    \item \relax \_total\_2
    \item \relax 2total
    \item \relax total 2
    \item \relax total-2
\end{enumerate}
\vspace{2mm}

\noindent\textbf{Q32.} \hfill {\small\color{LPUOrange}\textbf{[identifiers | Medium]}}\\[2pt]
Why does the statement \texttt{class = 5} raise an error in Python?

\begin{enumerate}[label=(\alph*),topsep=2pt,itemsep=1pt,partopsep=0pt,parsep=0pt]
    \item \relax \texttt{class} is a reserved keyword
    \item \relax integers cannot be assigned to names
    \item \relax \texttt{class} must be capitalized
    \item \relax the value must be a string
\end{enumerate}
\vspace{2mm}

\noindent\textbf{Q33.} \hfill {\small\color{LPUOrange}\textbf{[programming errors and debugging | Medium]}}\\[2pt]
A program runs without crashing but produces an incorrect average. What type of error is this?

\begin{enumerate}[label=(\alph*),topsep=2pt,itemsep=1pt,partopsep=0pt,parsep=0pt]
    \item \relax Runtime error
    \item \relax Logical error
    \item \relax Indentation error
    \item \relax Syntax error
\end{enumerate}
\vspace{2mm}

\noindent\textbf{Q34.} \hfill {\small\color{LPUOrange}\textbf{[programming errors and debugging | Medium]}}\\[2pt]
Dividing by zero in \texttt{x = 5 / 0} produces which kind of error?

\begin{enumerate}[label=(\alph*),topsep=2pt,itemsep=1pt,partopsep=0pt,parsep=0pt]
    \item \relax Runtime error
    \item \relax Syntax error
    \item \relax Semantic error
    \item \relax Logical error
\end{enumerate}
\vspace{2mm}

\noindent\textbf{Q35.} \hfill {\small\color{LPUOrange}\textbf{[expressions | Medium]}}\\[2pt]
What is the value of the expression \texttt{10 - 4 - 2} in Python?

\begin{enumerate}[label=(\alph*),topsep=2pt,itemsep=1pt,partopsep=0pt,parsep=0pt]
    \item \relax 0
    \item \relax 6
    \item \relax 8
    \item \relax 4
\end{enumerate}
\vspace{2mm}

\noindent\textbf{Q36.} \hfill {\small\color{LPUOrange}\textbf{[assignment statements | Medium]}}\\[2pt]
What happens when the following runs: \texttt{x = 5} then \texttt{x = x + "3"}?

\begin{enumerate}[label=(\alph*),topsep=2pt,itemsep=1pt,partopsep=0pt,parsep=0pt]
    \item \relax x becomes 8
    \item \relax TypeError is raised
    \item \relax x becomes "53"
    \item \relax x becomes 53
\end{enumerate}
\vspace{2mm}

\noindent\textbf{Q37.} \hfill {\small\color{LPUOrange}\textbf{[named constant | Medium]}}\\[2pt]
By convention, how should a named constant like the value of Pi be written in Python?

\begin{enumerate}[label=(\alph*),topsep=2pt,itemsep=1pt,partopsep=0pt,parsep=0pt]
    \item \relax pi
    \item \relax Pi
    \item \relax PI
    \item \relax *pi*
\end{enumerate}
\vspace{2mm}

\noindent\textbf{Q38.} \hfill {\small\color{LPUOrange}\textbf{[variables | Medium]}}\\[2pt]
After \texttt{a = 3} and \texttt{b = a} and then \texttt{a = 7}, what is the value of \texttt{b}?

\begin{enumerate}[label=(\alph*),topsep=2pt,itemsep=1pt,partopsep=0pt,parsep=0pt]
    \item \relax 3
    \item \relax 7
    \item \relax 10
    \item \relax None
\end{enumerate}
\vspace{2mm}

\noindent\textbf{Q39.} \hfill {\small\color{LPUOrange}\textbf{[need for a programming language | Medium]}}\\[2pt]
Why do we use high-level programming languages instead of machine code directly?

\begin{enumerate}[label=(\alph*),topsep=2pt,itemsep=1pt,partopsep=0pt,parsep=0pt]
    \item \relax They are easier for humans to read and write
    \item \relax They run faster than machine code
    \item \relax They require no translation to execute
    \item \relax They use only binary digits
\end{enumerate}
\vspace{2mm}

\noindent\textbf{Q40.} \hfill {\small\color{LPUOrange}\textbf{[introduction to python as a programming language | Medium]}}\\[2pt]
Which statement best describes Python as a programming language?

\begin{enumerate}[label=(\alph*),topsep=2pt,itemsep=1pt,partopsep=0pt,parsep=0pt]
    \item \relax It is an interpreted, dynamically typed language
    \item \relax It requires variable types to be declared
    \item \relax It converts code directly to machine code
    \item \relax It is a compiled, statically typed language
\end{enumerate}
\vspace{2mm}

\noindent\textbf{Q41.} \hfill {\small\color{LPUOrange}\textbf{[operator precedence and associativity | Hard]}}\\[2pt]
What is the result of evaluating \texttt{2 ** 3 ** 2} in Python?

\begin{enumerate}[label=(\alph*),topsep=2pt,itemsep=1pt,partopsep=0pt,parsep=0pt]
    \item \relax 81
    \item \relax 36
    \item \relax 512
    \item \relax 64
\end{enumerate}
\vspace{2mm}

\noindent\textbf{Q42.} \hfill {\small\color{LPUOrange}\textbf{[numeric data types | Hard]}}\\[2pt]
What does \texttt{0.1 + 0.2 == 0.3} evaluate to in Python, and why?

\begin{enumerate}[label=(\alph*),topsep=2pt,itemsep=1pt,partopsep=0pt,parsep=0pt]
    \item \relax \texttt{False}, due to floating-point representation error
    \item \relax \texttt{True}, since decimals are exact
    \item \relax \texttt{True}, after automatic rounding
    \item \relax \texttt{False}, because \texttt{==} is invalid for floats
\end{enumerate}
\vspace{2mm}

\noindent\textbf{Q43.} \hfill {\small\color{LPUOrange}\textbf{[type conversion and rounding | Hard]}}\\[2pt]
What is the output of \texttt{round(2.5) + round(3.5)}?

\begin{enumerate}[label=(\alph*),topsep=2pt,itemsep=1pt,partopsep=0pt,parsep=0pt]
    \item \relax 5
    \item \relax 8
    \item \relax 7
    \item \relax 6
\end{enumerate}
\vspace{2mm}

\noindent\textbf{Q44.} \hfill {\small\color{LPUOrange}\textbf{[simultaneous assignment | Hard]}}\\[2pt]
After executing \texttt{a, b = 1, 2} then \texttt{a, b = b, a + b}, what are the values of \texttt{a} and \texttt{b}?

\begin{enumerate}[label=(\alph*),topsep=2pt,itemsep=1pt,partopsep=0pt,parsep=0pt]
    \item \relax \texttt{a = 3, b = 2}
    \item \relax \texttt{a = 1, b = 3}
    \item \relax \texttt{a = 2, b = 3}
    \item \relax \texttt{a = 2, b = 4}
\end{enumerate}
\vspace{2mm}

\noindent\textbf{Q45.} \hfill {\small\color{LPUOrange}\textbf{[boolean types | Hard]}}\\[2pt]
What is the result of \texttt{True + True + False} in Python?

\begin{enumerate}[label=(\alph*),topsep=2pt,itemsep=1pt,partopsep=0pt,parsep=0pt]
    \item \relax 1
    \item \relax \texttt{TypeError}
    \item \relax 2
    \item \relax \texttt{True}
\end{enumerate}
\vspace{2mm}

\noindent\textbf{Q46.} \hfill {\small\color{LPUOrange}\textbf{[operators | Hard]}}\\[2pt]
What is the value of \texttt{-7 // 2} in Python?

\begin{enumerate}[label=(\alph*),topsep=2pt,itemsep=1pt,partopsep=0pt,parsep=0pt]
    \item \relax -4
    \item \relax -3
    \item \relax -3.5
    \item \relax 3
\end{enumerate}
\vspace{2mm}

\noindent\textbf{Q47.} \hfill {\small\color{LPUOrange}\textbf{[operators | Hard]}}\\[2pt]
What does \texttt{-7 \% 3} evaluate to in Python?

\begin{enumerate}[label=(\alph*),topsep=2pt,itemsep=1pt,partopsep=0pt,parsep=0pt]
    \item \relax -2
    \item \relax 2
    \item \relax 1
    \item \relax -1
\end{enumerate}
\vspace{2mm}

\noindent\textbf{Q48.} \hfill {\small\color{LPUOrange}\textbf{[expressions | Hard]}}\\[2pt]
What is the output of \texttt{2 + 3 * 4 ** 2 // 5 - 1}?

\begin{enumerate}[label=(\alph*),topsep=2pt,itemsep=1pt,partopsep=0pt,parsep=0pt]
    \item \relax 11
    \item \relax 10
    \item \relax 9
    \item \relax 8
\end{enumerate}
\vspace{2mm}

\noindent\textbf{Q49.} \hfill {\small\color{LPUOrange}\textbf{[augmented assignment operators | Hard]}}\\[2pt]
Given \texttt{x = 5}, what is the value of \texttt{x} after \texttt{x //= 2; x **= 2; x -= 1}?

\begin{enumerate}[label=(\alph*),topsep=2pt,itemsep=1pt,partopsep=0pt,parsep=0pt]
    \item \relax 5
    \item \relax 8
    \item \relax 3
    \item \relax 4
\end{enumerate}
\vspace{2mm}

\noindent\textbf{Q50.} \hfill {\small\color{LPUOrange}\textbf{[identifiers | Hard]}}\\[2pt]
Which of the following is a valid Python identifier?

\begin{enumerate}[label=(\alph*),topsep=2pt,itemsep=1pt,partopsep=0pt,parsep=0pt]
    \item \relax \texttt{class}
    \item \relax \texttt{total-sum}
    \item \relax \texttt{\_total2}
    \item \relax \texttt{2total}
\end{enumerate}
\vspace{2mm}

\noindent\textbf{Q51.} \hfill {\small\color{LPUOrange}\textbf{[type conversion and rounding | Hard]}}\\[2pt]
What is the output of \texttt{int("10", 2) + int("10")}?

\begin{enumerate}[label=(\alph*),topsep=2pt,itemsep=1pt,partopsep=0pt,parsep=0pt]
    \item \relax 22
    \item \relax 4
    \item \relax 20
    \item \relax 12
\end{enumerate}
\vspace{2mm}

\noindent\textbf{Q52.} \hfill {\small\color{LPUOrange}\textbf{[numeric data types | Hard]}}\\[2pt]
What is the type and value of \texttt{5 / 5} in Python 3?

\begin{enumerate}[label=(\alph*),topsep=2pt,itemsep=1pt,partopsep=0pt,parsep=0pt]
    \item \relax \texttt{float}, \texttt{1}
    \item \relax \texttt{int}, \texttt{1.0}
    \item \relax \texttt{float}, \texttt{1.0}
    \item \relax \texttt{int}, \texttt{1}
\end{enumerate}
\vspace{2mm}

\noindent\textbf{Q53.} \hfill {\small\color{LPUOrange}\textbf{[programming errors and debugging | Hard]}}\\[2pt]
A program runs without crashing but produces incorrect output. What type of error is this?

\begin{enumerate}[label=(\alph*),topsep=2pt,itemsep=1pt,partopsep=0pt,parsep=0pt]
    \item \relax Logical (semantic) error
    \item \relax Syntax error
    \item \relax Runtime exception
    \item \relax Indentation error
\end{enumerate}
\vspace{2mm}

\noindent\textbf{Q54.} \hfill {\small\color{LPUOrange}\textbf{[expressions | Hard]}}\\[2pt]
What is the result of \texttt{not True == False} in Python?

\begin{enumerate}[label=(\alph*),topsep=2pt,itemsep=1pt,partopsep=0pt,parsep=0pt]
    \item \relax \texttt{False}
    \item \relax \texttt{None}
    \item \relax \texttt{TypeError}
    \item \relax \texttt{True}
\end{enumerate}
\vspace{2mm}

\noindent\textbf{Q55.} \hfill {\small\color{LPUOrange}\textbf{[operators | Hard]}}\\[2pt]
What does \texttt{5 \& 3 | 2} evaluate to?

\begin{enumerate}[label=(\alph*),topsep=2pt,itemsep=1pt,partopsep=0pt,parsep=0pt]
    \item \relax 1
    \item \relax 7
    \item \relax 3
    \item \relax 2
\end{enumerate}
\vspace{2mm}

\noindent\textbf{Q56.} \hfill {\small\color{LPUOrange}\textbf{[assignment statements | Hard]}}\\[2pt]
After \texttt{x = y = [1, 2]; x.append(3)}, what is the value of \texttt{y}?

\begin{enumerate}[label=(\alph*),topsep=2pt,itemsep=1pt,partopsep=0pt,parsep=0pt]
    \item \relax \texttt{[1, 2, 3]}
    \item \relax \texttt{[3]}
    \item \relax \texttt{NameError}
    \item \relax \texttt{[1, 2]}
\end{enumerate}
\vspace{2mm}

\noindent\textbf{Q57.} \hfill {\small\color{LPUOrange}\textbf{[named constant | Hard]}}\\[2pt]
What is true about a named constant like \texttt{PI = 3.14159} in Python?

\begin{enumerate}[label=(\alph*),topsep=2pt,itemsep=1pt,partopsep=0pt,parsep=0pt]
    \item \relax It must be declared with a \texttt{const} keyword
    \item \relax It is stored in read-only memory automatically
    \item \relax It cannot be reassigned once defined
    \item \relax It is a convention only; Python does not enforce immutability
\end{enumerate}
\vspace{2mm}

\noindent\textbf{Q58.} \hfill {\small\color{LPUOrange}\textbf{[type conversion and rounding | Hard]}}\\[2pt]
What is the output of \texttt{round(3.14159, 2)} and its type?

\begin{enumerate}[label=(\alph*),topsep=2pt,itemsep=1pt,partopsep=0pt,parsep=0pt]
    \item \relax \texttt{3.14}, \texttt{str}
    \item \relax \texttt{3.14}, \texttt{float}
    \item \relax \texttt{3.15}, \texttt{float}
    \item \relax \texttt{3}, \texttt{int}
\end{enumerate}
\vspace{2mm}

\noindent\textbf{Q59.} \hfill {\small\color{LPUOrange}\textbf{[boolean types | Hard]}}\\[2pt]
What does the expression \texttt{0 or "" or [] or "Python"} return?

\begin{enumerate}[label=(\alph*),topsep=2pt,itemsep=1pt,partopsep=0pt,parsep=0pt]
    \item \relax \texttt{True}
    \item \relax \texttt{[]}
    \item \relax \texttt{"Python"}
    \item \relax \texttt{0}
\end{enumerate}
\vspace{2mm}

\noindent\textbf{Q60.} \hfill {\small\color{LPUOrange}\textbf{[introduction to python as a programming language | Hard]}}\\[2pt]
Which statement best describes Python's execution model?

\begin{enumerate}[label=(\alph*),topsep=2pt,itemsep=1pt,partopsep=0pt,parsep=0pt]
    \item \relax Source code is compiled to bytecode, then run by an interpreter (virtual machine)
    \item \relax Source code is translated to assembly and linked into an executable
    \item \relax Source code is compiled directly to native machine code before running
    \item \relax Source code is executed line-by-line with no intermediate representation
\end{enumerate}
\vspace{2mm}


\newpage
\section{Section B: Comprehensive Subjective Problem with Detailed Analytical Rubric}

\begin{subjectiveq}{1}{1}
Design and implement a robust Python algorithm to compute compound interest and loan amortization with exact rounding and error handling:
\begin{enumerate}[label=(\alph*)]
    \item Accept principal amount $P$, annual interest rate $r$ (percentage), loan tenure $t$ (years), and compounding frequency $n$ per year.
    \item Compute future compound value $A = P \left(1 + \frac{r}{100n}\right)^{nt}$ and effective annual rate $\text{EAR} = \left(1 + \frac{r}{100n}\right)^n - 1$.
    \item Validate inputs against type errors and domain boundaries ($P > 0, r > 0, t > 0, n \ge 1$).
    \item Format and round outputs strictly to 2 decimal places using Python string formatting.
\end{enumerate}
\end{subjectiveq}

\begin{solbox}
\textbf{Algorithmic Formulation and Implementation:}
\begin{pycode}
def calculate\_loan\_metrics(principal, rate, years, freq):\\
\quad if not (principal > 0 and rate > 0 and years > 0 and freq >= 1):\\
\quad \quad raise ValueError("Parameters must strictly satisfy P>0, r>0, t>0, n>=1")\\
~\\
\quad r\_dec = rate / 100.0\\
\quad amount = principal * ((1.0 + r\_dec / freq) ** (freq * years))\\
\quad ear = ((1.0 + r\_dec / freq) ** freq) - 1.0\\
\quad compound\_interest = amount - principal\\
~\\
\quad return \{\\
\quad \quad "maturity\_amount": round(amount, 2),\\
\quad \quad "total\_interest": round(compound\_interest, 2),\\
\quad \quad "effective\_rate\_pct": round(ear * 100.0, 4)\\
\quad \}\\
~\\
\# Test execution:\\
res = calculate\_loan\_metrics(100000, 8.5, 5, 12)\\
print(f"Maturity: \${res['maturity\_amount']:.2f}, Effective Rate: \{res['effective\_rate\_pct']\}\%")
\end{pycode}
\textbf{Complexity and Edge Analysis:} The arithmetic operations execute in $\mathcal{O}(1)$ time with $\mathcal{O}(1)$ auxiliary memory. Floor division and simultaneous type checks prevent zero division or negative base floating-point artifacts.
\end{solbox}

\begin{examinertip}
\begin{itemize}[leftmargin=1.5em, itemsep=2pt]
    \item \textbf{Algorithmic Formulation (3 Marks):} Correct problem decomposition, function signatures, and parameter documentation.
    \item \textbf{Syntactic Rigor \& Implementation (4 Marks):} Valid Python syntax, proper error handling, defensive input validation, and avoidance of anti-patterns.
    \item \textbf{Complexity \& Edge Analysis (3 Marks):} Accurate time and space asymptotic analysis $\mathcal{O}(n)$ and testing edge conditions (null inputs, zero division, empty collections).
\end{itemize}
\end{examinertip}
